% =====================================================================
%  People of the Holocaust -- Technology Evaluation
%
%  Formatted to match the project's Software Requirements Specification:
%  title page, table of contents, numbered sections, bordered tables.
%  Compiles with pdflatex on Overleaf.
%
%  NOTE: the table of contents needs two compilation passes to fill in.
%  Overleaf does this automatically; a single local pdflatex run leaves
%  the Contents page blank.
% =====================================================================

\documentclass[11pt]{article}

\usepackage[letterpaper,margin=1.25in]{geometry}
\usepackage{array}
\usepackage{longtable}
\usepackage[hidelinks]{hyperref}

\renewcommand{\arraystretch}{1.25}
\setlength{\parskip}{0.6em}
\setlength{\parindent}{0pt}

% column type for wrapping table cells
\newcolumntype{L}[1]{>{\raggedright\arraybackslash}p{#1}}

\begin{document}

% =====================================================================
%  TITLE PAGE
% =====================================================================
\begin{titlepage}
\begin{center}

\vspace*{1.2in}

{\LARGE\bfseries People of the Holocaust}\\[0.9em]
{\Large\bfseries Technology Evaluation}

\vfill

COSC 4920 - Principles of Design\\[0.6em]
Marquette University\\[1.2em]
Prepared by:\\[0.6em]
Version 1.0\\[0.6em]
\today

\vspace*{0.8in}

\end{center}
\end{titlepage}

% =====================================================================
%  CONTENTS
% =====================================================================
\tableofcontents
\newpage

% =====================================================================
\section{Revision History}

\begin{tabular}{|L{1in}|L{1.2in}|L{1.4in}|L{1.9in}|}
\hline
\textbf{Version} & \textbf{Date} & \textbf{Author} & \textbf{Description} \\
\hline
1.0 & \today & People of the Holocaust & Initial technology evaluation. \\
\hline
\end{tabular}

% =====================================================================
\section{Introduction}

\subsection{Purpose}

This document records the tools our team evaluated for the \textit{People of
the Holocaust} senior design project, and the reasoning behind each
decision. Each tool is assessed against the same three questions: what it
does, why it is useful, and whether it works for this project.

\subsection{Scope}

Three tools are covered: Supabase, MariaDB, and NLTK. Tools we chose not to
adopt are included alongside those we did, so that the reasoning behind each
decision remains available to whoever continues the project. We begin with
the tool we chose not to adopt, so that the constraint driving our database
decision is clear before we describe what we selected instead.

\subsection{Summary of Evaluations}

\begin{tabular}{|L{1.1in}|L{2.6in}|L{1.4in}|}
\hline
\textbf{Tool} & \textbf{Role considered for} & \textbf{Decision} \\
\hline
Supabase & Backend-as-a-service and hosted database & Not adopted \\
\hline
MariaDB & Relational database for parsed encyclopedia data & Adopted \\
\hline
NLTK & Tokenization and text parsing & Adopted \\
\hline
\end{tabular}

% =====================================================================
\section{Supabase}

\subsection{What it does}

Supabase is a platform for developers that sits on top of Postgres
databases. It works well as a backend service because it automatically turns
tables into REST APIs. It ships with the option to have row level security
(RLS), which needs to be manually configured. It also has authentication,
with options including one-time passcodes, email and password logins, and
anonymous sign-ins. Its free tier is plenty for most projects, but pauses
after seven days of inactivity.

\subsection{Why it is useful}

Supabase can be useful for many projects, since it sits on top of a Postgres
database, automatically turns tables into REST APIs, and ships all of this
in the free tier. One of the best things about Supabase is that, because it
is a backend-as-a-service, developers do not need to write a backend
completely from scratch and can simply use the APIs.

It fits particularly well in three situations. The first is a project that
needs a working backend quickly, without standing up and maintaining a
server; the previous team on this project chose Supabase for exactly this
reason, as it was quick and convenient and meant they never had to request a
Marquette server. The second is a project with user accounts, since
Supabase's authentication covers one-time passcodes, email and password
logins, and anonymous sign-ins without any of it being written by hand. The
third is prototypes and short-lived builds, such as a hackathon entry or a
one-semester project, where the limits of the free tier will never be
reached.

\subsection{Why it does not work for our project}

The previous team on this project ran on Supabase. Their interpreter signed
in with a stored service account and inserted parsed activity rows directly
into a Supabase table, while the intermediate people, place, and article
link tables were written out as CSV files. That arrangement was quick and
convenient for them, and it meant they never had to request a Marquette
server.

It does not hold up for what we need. Supabase, on the free tier, pauses
after seven days of inactivity. Given that our project needs to always be
displaying correct information, a database that can pause if nobody uses it
for a week is not viable. We also expect more traffic than the free tier is
designed to absorb once the site is in active use.

Supabase's authentication is one of its strongest features, but it is not a
feature we need. Every visitor to our site sees the same public information,
so there are no user accounts to manage and nothing for that part of the
platform to do.

We therefore moved to MariaDB, hosted on a Marquette Computer Science
server. Beyond removing the pausing problem, this makes the project clearer
to update and maintain, and it allows for an easier handoff than
transferring a Supabase project to a new team would be; instead, the new
team members are simply added to the server.

% =====================================================================
\section{MariaDB}

\subsection{What it does}

MariaDB is an open source relational database, forked from MySQL before
Oracle privatized it, known for being reliable and strongly supported within
the community. It uses standard SQL, with APIs for a number of languages
including Python, TypeScript, and JavaScript.

For this project, it stores our data on the MU server. Since we are scraping
and parsing data from the encyclopedia, we need a place to keep our parsed
data in tables so that our website knows where to find information.

\subsection{Why it is useful}

In general terms, MariaDB is lightweight and carries low overhead, it keeps
data local and therefore private, and it acts as a data engine for websites.
Each of those properties matters to this project in a specific way.

Because it is open source, we do not have to pay for usage. That matters
here, since the goal is for the site to receive more traffic once it is in
use, and a cost tied to usage would scale against us rather than with us.
Keeping the data local also means we do not have to share it with an
external third-party service: the scraped text is parsed locally on our own
machines, and the results go straight to our server and our website, so
nothing passes through a third party at any point. Day to day, creating
tables and querying data from them is straightforward, which keeps the
ongoing work of adding to and maintaining the database low.

\subsection{Why it works for our project}

As described above, MariaDB is an overall good open source, free relational
database management system. This lets us properly set up tables that have
relations with each other, and to get the information all we need to do is
query it with standard SQL, which means it is easier to pass on as well.

While there are other open source and licensed options for databases,
MariaDB being open source with no license needed provides another step of
simplicity and practicality. We want our database to last and to be easily
manageable, so this was the clear choice for us.

% =====================================================================
\section{NLTK}

\subsection{What it does}

NLTK is a Python library, short for Natural Language Toolkit. It has a
variety of built-in methods that make natural language processing easier.
The subsection of methods most relevant to this project are the tokenizer
methods. These are designed to take blocks of text and break them into
smaller tokens that can be analyzed individually and subsequently
reconstructed. Tokens can take the form of paragraphs, sentences, words, and
parts of words. The library also has methods that can gather tokens and
assign properties to larger tokens if smaller tokens also possess certain
properties.

\subsection{Why it is useful}

Analyzing text is a difficult process, as parsing through large strings
accurately is often tedious and time complex if numerous string comparisons
are being performed on every single character within a body of text.
Tokenizers reduce the comparisons being performed and make the process
simpler by splitting strings into smaller, more manageable pieces. They also
allow for the creation of tree structures where properties can be assigned
to blocks of text simply by traversing down those trees and gathering
properties of downstream nodes.

\subsection{Why it works for our project}

The \textit{People of the Holocaust} project involves taking large bodies of
text and attempting to extract relational information. For the purposes of
our project, we only consider information to be related if it appears in
close proximity or identically across an entire body of text.

For example, if we want to know who closely interacted with whom, we are
assuming they will appear in close proximity to each other in a body of
text, and so we would relate everyone that appeared in a paragraph as being
associated with each other. At the same time, if a name appears multiple
times across the encyclopedia, we want all mentions of one person to appear
on a single page. With a tokenizer you can simply see where identical person
tokens appear in the tree, as opposed to creating a needlessly complex
string comparison method.

\end{document}
